\section*{Annexe}
\subsection*{Implementations}

\subsubsection*{\underline{Displacement}}

\begin{verbatim}
Project
\   src
|   \   displacement_matrices
|   |   |   displacement_matrices.c
|   |   |   displacement_matrices.h
\end{verbatim}

\paragraph{gr\_mat\_displacement(gr\_mat\_t D, gr\_mat\_t A, disp\_type\_t type, gr\_ctx\_t ctx)}
Is the function taking as parameter the matrix $A$ and returns into the matrix $D$, indicating destination, the displacement matrix of $A$.
According to the type of displacement we pass via the parameter \texttt{type}, it calculates the implementations \eqref{imp:phi_plus} or \eqref{imp:phi_minus}.

\paragraph{int gr\_mat\_G\_H(gr\_mat\_t G, gr\_mat\_t H, gr\_mat\_t A, disp\_type\_t type, gr\_ctx\_t ctx))}
Is the function that computes the generator matrices $G$ and $H^T$ for an input matrix $A$.
Based on the displacement equations discussed previously, the equation \eqref{eq:SigmaLU}, the function first calculates the displacement matrix according to the specified type of displacement.
Since the displacement of a structured matrix (like a Toeplitz matrix) inherently results in a matrix with a low rank ($\alpha$),
$D$ can be factored into two smaller matrices such that $D = G H^T$.
The function achieves this by performing an LU decomposition on the displacement matrix $D$. It then extracts the independent columns to form $G$ and the corresponding rows to form $H^T$.
This is where the storage economy for structured matrices come from.

\paragraph{int gr\_mat\_reconstruct\_A(gr\_mat\_t A, gr\_mat\_t G, gr\_mat\_t H, disp\_type\_t type, gr\_ctx\_t ctx)}
\todo{RAFAEL -> V2 VERSION}


\subsubsection*{\underline{Compression}}\label{annexe:compression}

\begin{verbatim}
Project
\   src
|   \   compression
|   |   |   compression.c
|   |   |   compression.h
\end{verbatim}

\paragraph{int gr\_mat\_generator\_compress(gr\_mat\_t G\_d, gr\_mat\_t H\_d, gr\_ctx\_t ctx);} Is our function that implements the algorithms talked in chapter \ref{ch:compression}.